% !TeX program = pdflatex
% Una transformada y una dualidad: un núcleo de la teoría musical verificado por máquina en
%   Lean 4, de los conjuntos de alturas a los retículos de afinación.  -- versión ESPAÑOLA del
%   CORPUS PAPER (paraguas de las Notas #1/#2/#3). Traducción natural (no literal) de
%   corpus_paper.tex, en la línea de casa de los _es de las notas. Build: pdflatex (dos pasadas).

\documentclass[11pt]{article}
\usepackage[T1]{fontenc}
\usepackage[spanish,es-noindentfirst,es-noshorthands]{babel} % es-noshorthands: evita que `<>"` activos del español choquen con TikZ (>=...)
\usepackage{lmodern}
\usepackage{amsmath,amssymb,amsthm,booktabs,graphicx}
\usepackage{longtable}
\usepackage[hidelinks]{hyperref}
\usepackage{xurl}
\usepackage[table]{xcolor}
\usepackage[margin=1in]{geometry}
\usepackage{microtype}
\usepackage{float}
\usepackage[section]{placeins}
\usepackage{tikz}
\usetikzlibrary{arrows.meta,positioning,calc}

\definecolor{ctA}{HTML}{1B6F8C} % el hilo de la autocorrelación / DFT
\definecolor{ctB}{HTML}{B5530F} % el hilo de la dualidad
\setlength{\emergencystretch}{2em}

% --- colocación de floats: que LaTeX empaquete denso y no deje blanco perdido ---
\renewcommand{\topfraction}{0.92}
\renewcommand{\bottomfraction}{0.85}
\renewcommand{\textfraction}{0.07}
\renewcommand{\floatpagefraction}{0.78}
\setcounter{topnumber}{3}\setcounter{bottomnumber}{2}\setcounter{totalnumber}{5}

\newtheorem{theorem}{Teorema}[section]
\newtheorem{proposition}[theorem]{Proposición}
\newtheorem{lemma}[theorem]{Lema}
\newtheorem{remark}[theorem]{Observación}
\newtheorem{definition}[theorem]{Definición}

\newcommand{\lean}[1]{\texttt{#1}}
\newcommand{\Ztwelve}{\mathbb{Z}_{12}}
\newcommand{\Zn}[1]{\mathbb{Z}_{#1}}
\newcommand{\Dgrp}[1]{D_{#1}}
\newcommand{\Ahat}{\hat{A}}
\newcommand{\Fi}{\mathcal{F}^{-1}}
\newcommand{\Stab}{\mathrm{Stab}}
\newcommand{\question}[1]{\smallskip\noindent\textit{#1}\smallskip}
\newcommand{\audiobase}{https://github.com/karlesmarin/music-math/raw/main/media}
\newcommand{\audio}[1]{\href{\audiobase/#1}{\textcolor{ctA}{\footnotesize\,[$\flat$\,audio]}}}

\title{\bf Una transformada y una dualidad\\[5pt]
  \large\rm un núcleo de la teoría musical verificado por máquina en Lean~4,\\
  de los conjuntos de alturas a los retículos de afinación}
\author{Carles Mar\'in\thanks{Preparado con la asistencia de un sistema de IA (Claude, Anthropic).
La matemática es clásica; toda afirmación, delimitación de alcance y limitación es responsabilidad del
autor. La frontera de la contribución se enuncia con claridad en la Sección~\ref{sec:novelty}, y es el
núcleo de Lean~4 ---no el asistente--- lo que certifica las partes formalizadas.}\\[3pt]
  \normalsize Investigador independiente\\
  \normalsize \texttt{karlesmarin@gmail.com}}
\date{Junio de 2026\\[3pt]\normalsize DOI: \href{https://doi.org/10.5281/zenodo.20953768}{10.5281/zenodo.20953768}\quad\textbullet\quad
  \href{https://github.com/karlesmarin/music-math}{github.com/karlesmarin/music-math}}

\begin{document}
\maketitle

\begin{abstract}\noindent
La matemática de la teoría musical es matemática de verdad, pero casi toda vive en la prosa y en los
ejemplos trabajados. Aquí certificamos su núcleo en Lean~4 sobre Mathlib, en el universo cromático
$\Ztwelve$, y organizamos la biblioteca de modo que un solo objeto y un solo fenómeno hagan casi todo el
trabajo. El objeto es la autocorrelación, o espectro de potencia $|\Ahat|^2$, de un conjunto de alturas:
el vector de intervalos, el teorema del hexacordo de Babbitt, la homometría y la relación $Z$ leídas como
el problema cristalográfico de la fase, las escalas profundas, la rareza rítmica y los teoremas de notas
comunes son todos función suya, y la máxima uniformidad reconecta con él en la teoría de escalas. El
fenómeno es la dualidad: el grupo neo-riemanniano $PLR$ es el centralizador de la acción regular de $T/I$
sobre las $24$ tríadas (Crans--Fiore--Satyendra, vía el caso regular del teorema del centralizador de
Wielandt, un lema que añadimos sobre Mathlib), y la misma forma reaparece como una auto-dualidad de orden
$12$ centrada en el tritono de la clase de Petrushka $6$-$30$, y como la dualidad de retículos val--coma
de la teoría de temperamentos regulares. Cada teorema está sin \lean{sorry} y con auditoría de axiomas
limpia. La matemática de base es clásica. Lo que añadimos es la formalización ---que sepamos, la primera
exhaustiva en esta área, pues el trabajo previo en asistentes de demostración cubre sólo la acción
$T/I$---, junto con una lente organizadora de autocorrelación/DFT y dos resultados que no parecen haber
sido señalados en la literatura: la auto-dualidad de 6-30 y una taxonomía de fase unificada de los
invariantes de $|\Ahat|^2$.
\end{abstract}

\section{Introducción}\label{sec:intro}

\question{¿Qué puede certificar un asistente de demostración sobre la matemática de la música, y cuánto de
ella es en realidad una o dos ideas vistas desde muchos lados?}

La teoría de conjuntos de alturas, la teoría transformacional y neo-riemanniana, y la teoría de escalas y
afinación han construido, a lo largo de medio siglo, un cuerpo matemático sustancial sobre el grupo
cíclico $\Ztwelve$, el grupo diédrico, los retículos sobre $\mathbb{Z}$ y la transformada discreta de
Fourier: el teorema del hexacordo, la relación $Z$, la dualidad $T/I$--$PLR$, la máxima uniformidad, el
teorema de los tres huecos, la teoría de temperamentos regulares. Las demostraciones son correctas en el
sentido matemático corriente, pero se las \emph{cree} correctas, no se comprueban al nivel que exige un
núcleo de demostración; y la literatura arrastra una larga cola de casi-aciertos y de folclore enunciado
con más fuerza de la que se sostiene. Este trabajo certifica ese núcleo en Lean~4. Es el paraguas sobre
tres notas que aparecieron por separado ---sobre la auto-dualidad de 6-30~\cite{Note1}, la taxonomía de
fase~\cite{Note2} y el espectro del Tonnetz~\cite{Note3}--- y añade el tejido conectivo que las vuelve una
sola biblioteca.

\paragraph{Qué compra la verificación por máquina.} Nuestra motivación es un \emph{suelo de honestidad}.
Un enunciado verificado por máquina traza una línea dura entre lo demostrado y lo meramente plausible.
Donde un argumento clásico supone en silencio una no-degeneración, una formalización debe sacarla a la
superficie; donde un texto afirma algo de ``una escala'' pero sólo vale para la diatónica, el enunciado en
Lean nombra la instancia. Tratamos lo \textsc{no auditable} como información, no como fracaso: cuando sólo
se demuestra una instancia lo decimos, y nunca inflamos una primera formalización en una afirmación de
matemática nueva. La frontera de la contribución se traza en un solo sitio, con claridad, en la
Sección~\ref{sec:novelty}; todo lo demás es una primera formalización de un teorema publicado.

\paragraph{El modelo, el sustrato, el convenio de confianza.} Trabajamos sobre $\Ztwelve$. Un conjunto de
alturas es un \lean{Finset}\,$\Ztwelve$; una tríada es un punto de $\Ztwelve\times\mathrm{Bool}$; una
escala es una tupla ordenada cíclicamente. Nos apoyamos en dos piezas de Mathlib tal cual, en vez de
reconstruirlas: la transformada discreta de Fourier \lean{ZMod.dft} (Loeffler~\cite{Loeffler}), que nos da
$\Ahat$, el espectro de potencia, la inversión y Parseval gratis, y el grupo diédrico
\lean{DihedralGroup}\,$12$, que \emph{es} el grupo de transposiciones e inversiones del círculo cromático.
Cada teorema aquí está sin \lean{sorry}, y damos su auditoría de axiomas: la gran mayoría dependen sólo de
$[\mathrm{propext},\,\mathrm{Classical.choice},\,\mathrm{Quot.sound}]$, la base estándar de Mathlib, sin
\lean{sorryAx}; los pocos resultados de censo finito de la teoría de escalas invocan además
\lean{native\_decide} y por tanto cargan $\mathrm{Lean.ofReduceBool}$, que señalamos en cada caso mientras
mantenemos sin censo los motores analíticos que los rodean.

\paragraph{La biblioteca de un vistazo.} Para un lector que no trabaja en Lean, la
Tabla~\ref{tab:metrics} fija la escala de lo que aquí significa ``verificado por máquina''. Las fuentes
completas, los guiones testigo de Sage/GAP y las figuras están en
\url{https://github.com/karlesmarin/music-math}; cada teorema es comprobable con \lean{\#print axioms} y
cada guion es re-ejecutable.

\begin{table}[tbp]
\centering\small
\rowcolors{2}{ctA!7}{white}
\begin{tabular}{@{}l r@{}}
\toprule
\rowcolor{ctA!22}
\textbf{Métrica} & \textbf{Valor}\\
\midrule
Ficheros fuente de Lean (los tres pilares) & $9$\\
Teoremas y lemas & $275$\\
Declaraciones totales (incl.\ defs, instancias) & $372$\\
Líneas de Lean & ${\approx}\,3{,}800$\\
\lean{sorry} / \lean{admit} & $0$\\
Lemas nuevos aportados a la laguna de Mathlib & $1$ (Wielandt, caso regular)\\
Dependencia & Mathlib (Lean~4)\\
Base de axiomas & \texttt{propext}, \texttt{Classical.choice}, \texttt{Quot.sound}\\
\bottomrule
\end{tabular}
\caption{La biblioteca verificada por máquina detrás de los tres pilares. Cada teorema está sin
\lean{sorry}; las únicas salidas de la base de axiomas son los censos finitos por \lean{native\_decide} de
la~\S\ref{sec:p3-me}, que añaden \lean{Lean.ofReduceBool}.}
\label{tab:metrics}
\end{table}

\subsection{Dos ideas recurrentes}\label{sec:intro-threads}

Lo que hace de esto un solo trabajo, y no un catálogo, es que dos ideas organizan buena parte, y
reaparecen a través de los tres pilares (Figura~\ref{fig:arch}).

\question{Hilo~1. ¿Y si el vector de intervalos, el teorema del hexacordo, la homometría, las escalas
profundas, la rareza rítmica, las notas comunes y hasta la máxima uniformidad son todos el mismo objeto?}

\noindent\textbf{La autocorrelación / DFT.} Un solo objeto ---el espectro de potencia $|\Ahat|^2$,
equivalentemente la autocorrelación, equivalentemente la función de Patterson de los
cristalógrafos--- está detrás de una amplia familia de los invariantes del Pilar~1 (\S\ref{sec:pillar1}).
No son analogías: cada
uno es un teorema que iguala el invariante a una DFT inversa de $|\Ahat|^2$, de modo que la homometría es
\emph{exactamente} la pérdida de la fase y la transformada completa $\Ahat$ es exactamente la compleción
que falta. El mismo hilo reaparece en el Pilar~3: la escala maximalmente uniforme es aquella cuyo contenido
intervalar extremiza una energía convexa ---el maximizador del lado $|\Ahat|^2$---, de modo que la
singularidad de la diatónica es, otra vez, un enunciado sobre su autocorrelación.

\question{Hilo~2. Cuando un grupo actúa de forma simplemente transitiva, tiene un dual que es su propio
espejo; ¿qué le da eso a un acorde, a un acorde simétrico, a una afinación?}

\noindent\textbf{Dualidad: simplemente transitivo $\Rightarrow$ auto-centralizador.} El par dual de Lewin,
el grupo $T/I$ y el grupo neo-riemanniano $PLR$ ---cada uno el centralizador del otro sobre las $24$
tríadas--- es el teorema de Crans--Fiore--Satyendra del Pilar~2 (\S\ref{sec:pillar2}); su motor es el caso
regular del teorema del centralizador de Wielandt. La misma forma, sobre otro objeto, da la auto-dualidad
de orden $12$ centrada en el tritono del hexacordo $6$-$30$. Y rima una vez más en el Pilar~3, donde la
teoría de temperamentos regulares empareja el retículo de los vals con el de las comas sobre
$\mathbb{Z}$-módulos. Señalamos la rima pero no afirmamos un meta-teorema que unifique las dos
dualidades; viven en categorías distintas, y no tenemos tal teorema.

\begin{figure}[H]
\centering
\resizebox{0.8\textwidth}{!}{%
\begin{tikzpicture}[>={Stealth[length=2.4mm]},every node/.style={font=\footnotesize},
    box/.style={draw,rounded corners,inner sep=4pt,align=center,minimum height=8mm,text width=2.4cm},
    root/.style={draw,very thick,rounded corners,inner sep=4pt,align=center,minimum height=9mm,text width=2.7cm}]
  % ---- corona ----
  \node[draw,very thick,rounded corners,fill=black!5,font=\bfseries,inner sep=7pt] (crown) at (5.1,1.5)
       {La teoría musical en $\Ztwelve$};
  % ---- IZQUIERDA: la autocorrelación (Hilo 1) ----
  \node[root,ctA,fill=ctA!16] (pw) at (0.8,0) {$|\Ahat|^2$\\[-1pt]\scriptsize autocorrelación};
  \node[box] (iv)  at (-2.8,-1.8) {vector de intervalos};
  \node[box] (hom) at (0.8,-1.8)  {homometría / $Z$};
  \node[box] (ct)  at (4.4,-1.8)  {notas comunes,\\ escalas profundas, rareza};
  \node[box,ctA,fill=ctA!7] (me) at (-2.8,-3.5) {máxima uniformidad};
  \draw[ctA,->,thick] (pw)--(iv); \draw[ctA,->,thick] (pw)--(hom); \draw[ctA,->,thick] (pw)--(ct);
  \draw[ctA,->,thick] (iv)--(me);
  % ---- DERECHA: la dualidad (Hilo 2) ----
  \node[root,ctB,fill=ctB!16] (ti) at (9.6,0) {acción $T/I$\\[-1pt]\scriptsize regular};
  \node[box] (cen) at (9.6,-1.8) {centralizador};
  \node[box,ctB,fill=ctB!7] (plr) at (7.8,-3.5) {dualidad $PLR$};
  \node[box,ctB,fill=ctB!7] (s30) at (11.4,-3.5) {auto-dualidad de 6-30};
  \node[box,ctB,fill=ctB!7] (vc) at (11.4,-5.2) {dualidad val\,$\dashv$\,coma};
  \draw[ctB,->,thick] (ti)--(cen); \draw[ctB,->,thick] (cen)--(plr); \draw[ctB,->,thick] (cen)--(s30);
  \draw[ctB,->,thick] (s30)--(vc);
  % ---- conectores de la corona ----
  \draw[gray,dashed,->] (crown)--(pw); \draw[gray,dashed,->] (crown)--(ti);
\end{tikzpicture}}
\caption{La tesis en una imagen: el núcleo formalizado de la teoría musical en $\Ztwelve$ descansa sobre dos
mecanismos. \emph{Izquierda} ---la autocorrelación $|\Ahat|^2$ es la raíz común del vector de intervalos, la
homometría / la relación $Z$ y la familia de notas comunes (escalas profundas, rareza rítmica), y a través
del vector de intervalos alcanza la máxima uniformidad (Pilares~1 y~3). \emph{Derecha} ---una acción $T/I$
simplemente transitiva tiene un dual auto-centralizador, que se ramifica en la dualidad $PLR$ de las tríadas
y la auto-dualidad de 6-30, esta última continuando hasta la dualidad val\,$\dashv$\,coma de la teoría de
temperamentos (Pilares~2 y~3). Las dos columnas son estructuralmente paralelas; no afirmamos un teorema que
las una.}
\label{fig:arch}
\end{figure}

\paragraph{El recorrido.} El Pilar~1 (\S\ref{sec:pillar1}) desarrolla los invariantes de Fourier y cierra
en la taxonomía de fase. El Pilar~2 (\S\ref{sec:pillar2}) desarrolla la dualidad $T/I$--$PLR$ y la
auto-dualidad de 6-30. El Pilar~3 (\S\ref{sec:pillar3}) desarrolla escalas y afinación, donde los dos
hilos vuelven a encontrarse. La Sección~\ref{sec:novelty} enuncia con claridad qué es nuevo y qué es
primera formalización; la Sección~\ref{sec:frontier} recoge la frontera abierta ---lo que podemos
enunciar pero aún no certificar en el núcleo.

% ===========================================================================
% Los tres cuerpos de pilar, en la línea de casa (español).
% ===========================================================================
% pillar1_fourier_es.tex -- cuerpo del Pilar 1 (español), en la línea de casa.
% Sin preámbulo: usa las macros del maestro. La honestidad vive en la Sección \ref{sec:novelty}.

\section{Pilar 1: invariantes de Fourier de los conjuntos de alturas}\label{sec:pillar1}

\question{P\textsubscript{0}. El vector de intervalos, el teorema del hexacordo, la homometría, las
escalas profundas, la rareza rítmica, las notas comunes: ¿y si fueran todos un mismo objeto, visto
desde pliegues distintos?}

Codifica un conjunto de alturas $A\subseteq\Ztwelve$ por su indicadora $0/1$ y toma su transformada
discreta de Fourier $\Ahat=\mathcal{F}(\mathrm{ind}\,A)$, construida sobre la \lean{ZMod.dft} de Mathlib
(Loeffler~\cite{Loeffler}): $\Ahat(t)=\sum_{a\in A}e^{-2\pi i\,at/12}$ (\lean{Ahat}). Escribe
$\Ahat=|\Ahat|\,e^{i\varphi}$. La tesis de este pilar es que el \emph{espectro de potencia} $|\Ahat|^2$
(\lean{powerSpec}) ---equivalentemente la autocorrelación, equivalentemente la función de Patterson de los
cristalógrafos--- es un solo objeto detrás de una amplia familia de los invariantes clásicos, y que la
frontera de lo que puede ver es \emph{exactamente} la fase $\varphi$. Cada enunciado de abajo es un teorema en
\lean{Fourier.lean}/\lean{IntervalVector.lean} (Tabla~\ref{tab:p1}), sin \lean{sorry} y axioma-limpio.

\subsection{El vector de intervalos es la autocorrelación}\label{sec:p1-ivkey}

\question{P\textsubscript{1}. ¿Por qué es invariante el vector de intervalos bajo transposición e
inversión, y qué es, espectralmente?}

El vector de intervalos cuenta, para cada clase de intervalo $k$, los pares ordenados de $A$ con esa clase
de diferencia (\lean{IVraw}); es invariante bajo transposición (\lean{IVraw\_tpose}) y bajo inversión
(\lean{IVraw\_invert}) porque ambas sólo permutan diferencias. Su identidad espectral es el eje del pilar.
Con la autocorrelación $\mathrm{autocorr}(A,d)=\#\{(a,b)\in A^2:a-b=d\}$ (\lean{autocorr}), el teorema de
Wiener--Khinchin vale al pie de la letra, $\mathrm{autocorr}(A,d)=\Fi(|\Ahat|^2)(d)$
(\lean{autocorr\_eq\_invDFT}), y reagrupar por clase de intervalo da la pieza clave
\[
  \mathrm{IVraw}(A,k)=\textstyle\sum_{\mathrm{ivc}\,d=k}\Fi(|\Ahat|^2)(d)
  \qquad(\lean{IVraw\_eq\_invDFT\_power}).
\]
El vector de intervalos es, por tanto, función de $|\Ahat|^2$ \emph{solo}: no puede ver la fase. Todo lo de
la subsección siguiente cabalga sobre esto.

\subsection{La familia ciega a la fase: hexacordo, notas comunes, escalas profundas, rareza rítmica}
\label{sec:p1-family}

\question{P\textsubscript{2}. ¿Cuánto del piso clásico se desprende de ``es función de $|\Ahat|^2$''?}

\emph{El teorema del hexacordo de Babbitt.} Para $|A|=6$ y todo $t\neq0$, $A$ y su complemento tienen
transformadas opuestas (\lean{Ahat\_compl}), luego iguales espectros de potencia
(\lean{hexachord\_powerSpec\_eq}) y, por la pieza clave, igual contenido intervalar,
$\mathrm{IVraw}(A,k)=\mathrm{IVraw}(A^{\mathsf c},k)$ (\lean{hexachord\_IVraw\_eq}): un hexacordo y su
complemento tienen el mismo vector de intervalos (Babbitt; la prueba por DFT es de Amiot~\cite{Amiot}).

\emph{Notas comunes y el tritono.} Las alturas que quedan fijas bajo $T_k$ son exactamente la
autocorrelación, $\#(A\cap T_kA)=\Fi(|\Ahat|^2)(k)$ (\lean{commonTones\_eq\_invDFT}); sobre una clase de
intervalo $1\le k\le5$ la fibra es un par, así que $\mathrm{IVraw}(A,k)=2\,\#(A\cap T_kA)$, mientras que
sobre la clase del tritono $6$ es un singleton y el factor $2$ desaparece (\lean{IVraw\_tritone}). El
tritono $T_6$ es el mismo elemento central que organiza el Pilar~2.

\emph{Escalas profundas y rareza rítmica.} Una escala es \emph{profunda} cuando sus seis multiplicidades de
clase de intervalo son distintas (\lean{IsDeep}) ---por la identidad del tritono, exactamente cuando cada
transposición retiene un número distinto de notas (\lean{IsDeep\_iff\_commonTones\_nodup}); la diatónica es
profunda, la tríada aumentada no (\lean{diatonic\_isDeep}, \lean{augmented\_not\_isDeep}). Un ritmo cíclico
tiene la propiedad de rareza rítmica cuando no hay dos ataques a un tritono de distancia, es decir, cuando
la autocorrelación se anula en $6$ (\lean{rop\_iff\_autocorr\_zero}), lo que fuerza $|A|\le6$
(\lean{rop\_card\_le}).

\subsection{Homometría: la fase que el vector de intervalos no puede ver}\label{sec:p1-homometry}

\question{P\textsubscript{3}. Si el vector de intervalos es ciego a la fase, dos conjuntos pueden
compartirlo y aun así diferir. ¿Cuál es la información que falta?}

Dos conjuntos son \emph{homométricos} cuando comparten la autocorrelación (\lean{Homometric});
equivalentemente, tienen el mismo espectro de potencia (\lean{homometric\_iff\_powerSpec\_eq}),
equivalentemente la misma intensidad de difracción $|\Ahat|^2$ en cada frecuencia
(\lean{homometric\_iff\_normSq\_dft\_eq}). Esto es, al pie de la letra, el problema de la fase de la
cristalografía de rayos X: la autocorrelación es la función de \emph{Patterson} (\lean{Patterson};
Patterson 1934~\cite{Patterson}), $|\Ahat|^2$ la intensidad medida, las estructuras homométricas los
acordes $Z$-relacionados (Forte~\cite{Forte}; la correspondencia $Z\leftrightarrow$Patterson es de
Mandereau et al.~\cite{MAAA}). El testigo canónico es el par de tetracordes de todos los intervalos
(Figura~\ref{fig:p1-zpair}): $\mathrm{Homometric}\,\{0,1,4,6\}\,\{0,1,3,7\}$
(\lean{allIntervalTetrachords\_homometric}, por \lean{decide})\,\audio{z_pair.wav} ---el par 4-Z15 / 4-Z29
de Forte, misma magnitud, fase distinta, no $T/I$-relacionados. La información que falta es exactamente
$\varphi$.

\begin{figure}[ht]
\centering
\begin{tikzpicture}[scale=0.42,every node/.style={font=\footnotesize}]
  \begin{scope}[shift={(0,0)}]
    \draw[gray!55] (0,0) circle (3);
    \foreach \k in {0,...,11}{
      \pgfmathtruncatemacro{\ina}{(\k==0||\k==1||\k==4||\k==6)?1:0}
      \ifnum\ina=1 \filldraw[ctA] ({90-30*\k}:3) circle (0.18);
      \else \draw[fill=white] ({90-30*\k}:3) circle (0.14); \fi}
    \node at (0,-3.9) {$\{0,1,4,6\}$ = 4-Z15};
  \end{scope}
  \begin{scope}[shift={(9,0)}]
    \draw[gray!55] (0,0) circle (3);
    \foreach \k in {0,...,11}{
      \pgfmathtruncatemacro{\inb}{(\k==0||\k==1||\k==3||\k==7)?1:0}
      \ifnum\inb=1 \filldraw[ctB] ({90-30*\k}:3) circle (0.18);
      \else \draw[fill=white] ({90-30*\k}:3) circle (0.14); \fi}
    \node at (0,-3.9) {$\{0,1,3,7\}$ = 4-Z29};
  \end{scope}
  \node at (4.5,0) {\large $\overset{|\Ahat|^2}{=}$};
\end{tikzpicture}
\caption{El par $Z$-relacionado (homométrico) $\{0,1,4,6\}$ y $\{0,1,3,7\}$: idéntico espectro de potencia
$|\Ahat|^2=[16,2,4,2,4,2,4,2,4,2,4,2]$ e idéntico vector de intervalos, y sin embargo no relacionados por
ninguna transposición ni inversión. El vector de intervalos ---función de $|\Ahat|^2$ solo--- no los puede
distinguir; sólo la fase de $\Ahat$ puede.}
\label{fig:p1-zpair}
\end{figure}

\subsection{El vector de sumas: un primer atisbo de la fase}\label{sec:p1-sum}

\question{P\textsubscript{4}. ¿Puede algún invariante ver \emph{algo} de la fase que el vector de
intervalos pierde?}

Donde el vector de intervalos cuenta notas fijas bajo \emph{transposición} (el emparejamiento por
diferencia $a-b$, gobernado por $|\Ahat|^2$), el \emph{vector de sumas / índices} cuenta notas fijas bajo
\emph{inversión} $I_j:x\mapsto j-x$ (el emparejamiento por suma $a+b$): $\#(A\cap I_jA)=\mathrm{sumcorr}(A,j)$
(\lean{commonTonesInv\_eq\_sumcorr}), y esto es la DFT inversa del \emph{coeficiente} al cuadrado $\Ahat^2$,
no de la magnitud al cuadrado ---así que retiene fase (\lean{sumcorr\_eq\_invDFT}). Ya parte el par
homométrico: en torno al eje $j=0$, $\{0,1,4,6\}$\,\audio{inversion_4z15.wav} retiene dos notas y
$\{0,1,3,7\}$\,\audio{inversion_4z29.wav} retiene una (\lean{sumVector\_distinguishes\_homometric}, por
\lean{decide}). La fase que el vector de intervalos no ve, un conteo inversional la ve ---en parte.

\subsection{La cúspide: la taxonomía de fase}\label{sec:p1-capstone}

\question{P\textsubscript{5}. ¿Qué haría falta para ver \emph{toda} la fase?}

La transformada completa. $\Ahat$ es un invariante completo: $\Ahat_A=\Ahat_B\Rightarrow A=B$
(\lean{Ahat\_injective}), con bicondicional \lean{Ahat\_eq\_iff} ---la prueba es corta y clásica, pues
\lean{ZMod.dft} es una \lean{LinearEquiv}, luego inyectiva, y la indicadora recupera la pertenencia
(\lean{ind\_injective}). Pasar a una clase de conjuntos cocienta exactamente dos operaciones de fase: una
transposición es una rampa de fase lineal en frecuencia sobre $\Ahat$ (\lean{tpose\_iff\_Ahat\_ramp}) y una
inversión una conjugación-más-rampa (\lean{inv\_iff\_Ahat\_conj\_ramp}), de modo que $B$ está en la órbita
$T/I$ de $A$ syss $\Ahat_B$ iguala a $\Ahat_A$ salvo una rampa, o salvo conjugación-y-rampa
(\lean{setClass\_iff\_Ahat\_mod\_phase}). El pilar cierra, pues, como una escalera con dos fronteras
verificadas por máquina, una junto a otra: el piso de sólo magnitud $|\Ahat|^2$, donde la homometría
colapsa acordes distintos (P\textsubscript{3}), y el invariante completo $\Ahat$, donde nada se pierde
(P\textsubscript{5}). La fase es exactamente el dato que completa el vector de intervalos hasta un
invariante total; volvemos a lo que esto le dice a la auto-dualidad del Pilar~2 en
la~\S\ref{sec:p2-sixthirty}.

\begin{figure}[ht]
\centering
\resizebox{0.82\textwidth}{!}{%
\begin{tikzpicture}[every node/.style={font=\footnotesize},>={Stealth[length=2.4mm]}]
  \node[draw,rounded corners,fill=ctA!8,text width=8.6cm,align=center,minimum height=8mm] (l1)
    {$|\Ahat|^2$ \;---\; sólo magnitud (ciego a la fase): la homometría colapsa $4$-Z15 $\equiv$ $4$-Z29};
  \node[draw,rounded corners,fill=ctA!16,text width=8.6cm,align=center,minimum height=8mm,above=9mm of l1] (l2)
    {$\Ahat^2$ \;---\; vector de sumas / índices (fase parcial): \emph{parte} el par $Z$};
  \node[draw,rounded corners,ctB,thick,fill=ctB!10,text width=8.6cm,align=center,minimum height=8mm,above=9mm of l2] (l3)
    {$\Ahat$ \;---\; la transformada completa (toda la fase): un invariante \emph{completo}};
  \draw[->,thick] (l1) -- (l2) node[midway,right=1pt,font=\scriptsize]{$+$ fase parcial};
  \draw[->,thick] (l2) -- (l3) node[midway,right=1pt,font=\scriptsize]{$+$ el resto};
\end{tikzpicture}}
\caption{La taxonomía como una escalera de información creciente. El peldaño inferior, el espectro de
potencia $|\Ahat|^2$, es ciego a la fase, y la homometría colapsa sobre él el par $4$-Z15 / $4$-Z29
(\lean{homometric\_iff\_normSq\_dft\_eq}); el coeficiente al cuadrado $\Ahat^2$ restaura fase suficiente
para partir el par (\lean{sumVector\_distinguishes\_homometric}); la transformada completa $\Ahat$ es un
invariante completo (\lean{Ahat\_injective}). Las dos fronteras ---ciega a la magnitud y completa--- se
certifican una junto a otra.}
\label{fig:p1-ladder}
\end{figure}

\begin{table}[ht]
\centering\scriptsize\setlength{\tabcolsep}{4pt}
\rowcolors{2}{ctA!7}{white}
\begin{tabular}{@{}p{0.35\textwidth} p{0.57\textwidth}@{}}
\toprule
\rowcolor{ctA!22}
\textbf{Teorema} & \textbf{Enunciado}\\
\midrule
\lean{Ahat\_zero}, \lean{Ahat\_tpose}, \lean{Ahat\_compl} & $\Ahat(0)=|A|$; transposición $=$ rampa de fase; complemento $=$ negación fuera de DC.\\
\lean{IVraw\_tpose}, \lean{IVraw\_invert} & vector de intervalos invariante bajo $T$ e $I$.\\
\lean{autocorr\_eq\_invDFT} & Wiener--Khinchin: $\mathrm{autocorr}=\Fi(|\Ahat|^2)$.\\
\lean{IVraw\_eq\_invDFT\_power} & \textbf{pieza clave}: vector de intervalos $=\Fi(|\Ahat|^2)$ reagrupado por clase.\\
\lean{hexachord\_powerSpec\_eq}, \lean{hexachord\_IVraw\_eq} & Babbitt: hexacordo y complemento comparten $|\Ahat|^2$, luego el contenido intervalar.\\
\lean{commonTones\_eq\_invDFT}, \lean{IVraw\_tritone} & notas comunes $=\Fi(|\Ahat|^2)$; factor del tritono (fibra singleton).\\
\lean{IsDeep\_iff\_commonTones\_nodup} & profunda $\iff$ conteos de notas comunes por transposición distintos.\\
\lean{rop\_iff\_autocorr\_zero}, \lean{rop\_card\_le} & rareza rítmica $\iff \mathrm{autocorr}(\cdot,6)=0$; fuerza $|A|\le6$.\\
\lean{homometric\_iff\_normSq\_dft\_eq} & homometría $\iff$ igual $|\Ahat|^2$ (el problema discreto de la fase).\\
\lean{allIntervalTetrachords\_homometric} & $\{0,1,4,6\}\sim\{0,1,3,7\}$ (4-Z15 / 4-Z29), el testigo $Z$.\\
\lean{sumcorr\_eq\_invDFT}, \lean{sumVector\_distinguishes\_homometric} & vector de sumas $=\Fi(\Ahat^2)$ (fase parcial); parte el par $Z$.\\
\lean{Ahat\_injective}, \lean{Ahat\_eq\_iff} & la transformada completa $\Ahat$ es un invariante completo.\\
\lean{setClass\_iff\_Ahat\_mod\_phase} & clase de conjuntos $=\Ahat$ salvo una rampa de fase / conjugación-y-rampa.\\
\bottomrule
\end{tabular}
\caption{El Pilar~1 en \lean{Fourier.lean} / \lean{IntervalVector.lean}: los invariantes de la escalera de
la fase, sin \lean{sorry} y con auditoría de axiomas $[\mathrm{propext},\mathrm{Classical.choice},
\mathrm{Quot.sound}]$ (los testigos por \lean{decide} libres de elección). La novedad se adjudica en la
Sección~\ref{sec:novelty}.}
\label{tab:p1}
\end{table}
     % \S\ref{sec:pillar1}
% pillar2_duality_es.tex -- cuerpo del Pilar 2 (español), en la línea de casa.
% Sin preámbulo: usa las macros del maestro. La honestidad vive en la Sección \ref{sec:novelty}.

\section{Pilar 2: dualidad transformacional}\label{sec:pillar2}

\question{P\textsubscript{0}. Un grupo que actúa de forma simplemente transitiva tiene un \emph{dual} que
lo refleja. ¿Qué le da ese dual a una tríada, y qué le pasa ante un acorde \emph{simétrico}?}

La tradición neo-riemanniana lee los acordes no por sus notas sino por el grupo de operaciones que los
relaciona, y empareja cada grupo de transformaciones con un dual: dos subgrupos de $\mathrm{Sym}(\Omega)$
son \emph{duales} en el sentido de Lewin (\cite{Lewin}, p.~157) cuando cada uno es el centralizador del
otro y ambos actúan de forma simplemente transitiva. Este pilar verifica por máquina esa dualidad para las
$24$ tríadas, y luego la sigue al único sitio donde hace algo nuevo ---un acorde simétrico, donde el
tritono es el centro del grupo.

\subsection{La acción de $T/I$ sobre las tríadas y el grupo $PLR$}\label{sec:p2-ti}

Una \emph{tríada} es un par $(\mathrm{fundamental},\mathrm{calidad})\in\Ztwelve\times\mathrm{Bool}$ (mayor
$\{r,r{+}4,r{+}7\}$ o menor $\{r,r{+}3,r{+}7\}$); hay $24$. El grupo $T/I$ es el \lean{DihedralGroup}\,$12$
de Mathlib, que actúa sobre las tríadas de forma simplemente transitiva ---la representación regular por la
izquierda (\lean{card\_TIgrp}\,$=24$, con las instancias \lean{MulAction}, \lean{IsPretransitive} y
\lean{FaithfulSMul}). Las operaciones neo-riemannianas $P,L,R$ son las tres involuciones que intercambian
una tríada con la que comparte dos notas comunes (\lean{Pf},\lean{Lf},\lean{Rf}); cada una conmuta con toda
operación de $T/I$ (\lean{Pf\_comm}, \lean{Lf\_comm}, \lean{Rf\_comm}), y cada una es un movimiento de
conducción de voces mínimo: $P$ y $L$ mueven una voz un semitono (\lean{P\_size\_one},
\lean{L\_size\_one}), $R$ un tono entero (\lean{R\_size\_two}), reteniendo cada una las otras dos
notas (\lean{P\_holds\_two}, \lean{L\_holds\_two}, \lean{R\_holds\_two}) ---la parsimonia sobre la que Cohn
\cite{Cohn} construyó el Tonnetz.

\subsection{$PLR$ es el centralizador de $T/I$ (Crans--Fiore--Satyendra)}\label{sec:p2-cfs}

\question{P\textsubscript{1}. La tríada es asimétrica, así que su órbita tiene $24$ formas. ¿Es de verdad
el grupo $PLR$ el dual \emph{entero}, y por qué?}

La dualidad es el caso regular del teorema clásico del centralizador: el centralizador de un grupo de
permutaciones regular es su representación regular opuesta (Dixon--Mortimer~\cite{DM} \S4.2,
Cameron~\cite{Cam}). Su motor es un solo lema ---un elemento del centralizador que fija un punto es la
identidad (\lean{centralizing\_fixedPoint\_eq\_one})--- que \emph{no estaba} en Mathlib y que aportamos. De
él: $\langle P,L,R\rangle$ está en el centralizador (\lean{PLR\_le\_centralizer}), el centralizador tiene a
lo sumo $24$ elementos y $PLR$ tiene exactamente $24$ (\lean{card\_PLRgrp}), y ambos coinciden,
\[
  \mathrm{centralizador}(T/I)=\langle P,L,R\rangle\qquad(\lean{centralizer\_eq\_PLR}),
\]
la dualidad de Crans--Fiore--Satyendra~\cite{CFS}.

\subsection{La auto-dualidad de 6-30 centrada en el tritono}\label{sec:p2-sixthirty}

\question{P\textsubscript{2}. ¿Qué le pasa al dual cuando el objeto es un acorde \emph{simétrico}, cuya
órbita es menor que $24$?}

Una clase simétrica tiene estabilizador $T/I$ no trivial, así que su órbita es menor y la acción deja de
ser libre; un dual reducido simplemente transitivo sobrevive exactamente cuando el estabilizador es
\emph{normal} en $\Dgrp{12}$, y entre las simetrías de orden $2$ eso singulariza el centro $\{T_0,T_6\}$
---el tritono. La única clase de conjuntos no trivial en $\Ztwelve$ con justo esa simetría es la
\textbf{6-30} de Forte $=\{0,1,3,6,7,9\}$, la clase del acorde de Petrushka de Stravinsky
(Figura~\ref{fig:p2-clock}): una enumeración exhaustiva de las $222$ clases de conjuntos no triviales de
$\Ztwelve$ confirma que no hay otra clase cuyo estabilizador $T/I$ sea exactamente $\{T_0,T_6\}$ (una
comprobación finita, recogida en la Nota~\#1~\cite{Note1}). Es la unión de los tres pares centrales
$\{0,6\},\{1,7\},\{3,9\}$, así que la
fija el tritono $T_6$ (\lean{T6\_invariant})\,\audio{t6_symmetry.wav} y no tiene simetría inversional
(\lean{not\_inversionally\_symmetric}). Su órbita $T/I$ colapsa por tanto a $12$ (\lean{card\_O}); el
tritono cae en el núcleo de la acción sobre la órbita (\lean{Ker\_eq}\,$=\{1,r_6\}$); y el dual ---el
centralizador de la acción reducida--- tiene orden $12$ (\lean{card\_Cgrp}, construido como la
representación regular opuesta), es transitivo y por tanto regular (\lean{Cgrp\_transitive}), y lleva la
huella diédrica (\lean{dihedral\_fingerprint}: un generador $s$ de orden $6$, una involución $t\neq1$ con
$tst=s^{-1}$ y $st\neq ts$), fijándolo como diédrico de orden $12$. Así, las doce formas del acorde de
Petrushka forman un único espacio armónico navegable\,\audio{orbit_12.wav}, el análogo para acordes
simétricos de la red $P/L/R$ de las tríadas ---el tritono que Stravinsky buscó de oído \emph{es} la
simetría central $T_6$, y ese centro es exactamente lo que dota a 6-30 de su auto-dualidad de orden 12.

\begin{figure}[ht]
\centering
\begin{tikzpicture}[scale=1.2,every node/.style={font=\small}]
  \draw[thick] (0,0) circle (1);
  \foreach \a/\b in {0/6,1/7,3/9}{
    \draw[ctB,line width=1.3pt] ({90-30*\a}:1) -- ({90-30*\b}:1);}
  \foreach \k in {0,...,11}{
    \pgfmathtruncatemacro{\ina}{(\k==0||\k==1||\k==3||\k==6||\k==7||\k==9)?1:0}
    \ifnum\ina=1 \filldraw[black] ({90-30*\k}:1) circle (0.05);
      \node at ({90-30*\k}:1.2) {\textbf{\k}};
    \else \draw[fill=white] ({90-30*\k}:1) circle (0.045);
      \node[gray] at ({90-30*\k}:1.2) {\k}; \fi}
\end{tikzpicture}
\caption{La \textbf{6-30}$=\{0,1,3,6,7,9\}$ en el reloj de clases de altura: la unión de los tres pares
centrales $\{0,6\},\{1,7\},\{3,9\}$ (diámetros gruesos de tritono), por lo que la fija $T_6$ (una rotación
de $180^\circ$) y, sin simetría inversional, la estabiliza exactamente $\{T_0,T_6\}$. Ese estabilizador
central es lo que reduce la órbita a $12$ y vuelve diédrico de orden $12$ al dual.}
\label{fig:p2-clock}
\end{figure}

\begin{figure}[ht]
\centering
\begin{tikzpicture}[scale=0.85,every node/.style={font=\footnotesize},>={Stealth[length=2mm]}]
  \foreach \k in {0,...,5}{
    \coordinate (O\k) at ({90-60*\k}:2.4);
    \coordinate (I\k) at ({90-60*\k}:1.2);}
  \foreach \k in {0,...,5}{
    \pgfmathtruncatemacro{\nextk}{mod(\k+1,6)}
    \draw[ctA,->,thick] (O\k) -- (O\nextk);
    \draw[ctB,->,thick] (I\nextk) -- (I\k);
    \draw[gray,dashed] (O\k) -- (I\k);}
  \foreach \k in {0,...,5}{
    \filldraw[fill=ctA!12,draw=ctA,thick] (O\k) circle (0.32) node[black]{$T_{\k}A$};
    \filldraw[fill=white,draw=ctB,thick] (I\k) circle (0.32) node[black]{$I_{\k}A$};}
\end{tikzpicture}
\caption{El grupo dual de orden $12$, $\Dgrp{6}$, actuando de forma simplemente transitiva sobre las $12$
formas de $6$-$30$ ---su grafo de Cayley, el análogo para acordes simétricos de la red $P/L/R$ de las
tríadas. Los nodos exteriores son las seis formas de transposición $T_kA$ (con $T_kA=T_{k+6}A$
identificadas), los interiores las seis de inversión $I_kA$; las flechas sólidas son un generador de
rotación de orden $6$ (la orientación se invierte en la clase interior, como en todo grafo de Cayley
diédrico), las aristas discontinuas un generador de reflexión de orden $2$. Cualquier forma del acorde de
Petrushka alcanza cualquier otra por un movimiento dual único (\lean{Cgrp\_transitive}), así que las doce
formas se vuelven un único espacio armónico navegable.}
\label{fig:p2-cayley}
\end{figure}

\subsection{El hilo de la dualidad entre pilares}\label{sec:p2-thread}

La forma de aquí ---una acción simplemente transitiva y su espejo auto-centralizador--- vuelve en el
Pilar~3. Allí, la teoría de temperamentos regulares empareja el retículo de los vals con el de las comas
sobre $\mathbb{Z}$-módulos: la misma clase de dualidad, una estructura generada y su anulador.
Es un paralelismo estructural, no un teorema que unifique las dos (Sección~\ref{sec:novelty}). Y la
auto-dualidad de un acorde simétrico es, en términos de Fourier, un hecho de fase que $|\Ahat|$ no puede
ver (\S\ref{sec:p1-capstone}) ---la simetría $T_6$ es soporte en frecuencias pares, un enunciado de
magnitud, pero qué acordes simétricos son auto-duales necesita la fase, la misma ceguera que la relación
$Z$.

\begin{table}[ht]
\centering\scriptsize\setlength{\tabcolsep}{4pt}
\rowcolors{2}{ctB!7}{white}
\begin{tabular}{@{}p{0.37\textwidth} p{0.55\textwidth}@{}}
\toprule
\rowcolor{ctB!22}
\textbf{Teorema} & \textbf{Enunciado}\\
\midrule
\multicolumn{2}{@{}l}{\emph{el motor de dualidad, \lean{NeoRiemannian.lean}}}\\
\lean{card\_TIgrp} & $T/I$ actúa sobre las $24$ tríadas como la rep.\ regular de $\Dgrp{12}$.\\
\lean{Pf\_comm}, \lean{Lf\_comm}, \lean{Rf\_comm} & $P,L,R$ conmutan cada uno con la acción de $T/I$.\\
\lean{P\_size\_one}, \lean{L\_size\_one}, \lean{R\_size\_two} & tamaños de conducción de voces mínimos de $P,L,R$.\\
\lean{centralizing\_fixedPoint\_eq\_one} & caso regular de Wielandt (el centralizador es semirregular) --- laguna de Mathlib.\\
\lean{centralizer\_eq\_PLR} & $\mathrm{centralizador}(T/I)=\langle P,L,R\rangle$ (dualidad CFS); \lean{card\_PLRgrp}$=24$.\\
\multicolumn{2}{@{}l}{\emph{la auto-dualidad de 6-30, \lean{SixThirty.lean}}}\\
\lean{T6\_invariant}, \lean{not\_inversionally\_symmetric} & $A=\{0,1,3,6,7,9\}$ fija por $T_6$, sin simetría de inversión.\\
\lean{card\_O}, \lean{Ker\_eq} & órbita $12$; núcleo de la acción $=\{T_0,T_6\}=Z(\Dgrp{12})$.\\
\lean{card\_Cgrp}, \lean{Cgrp\_transitive} & dual (centralizador sobre la órbita) orden $12$, regular.\\
\lean{dihedral\_fingerprint} & orden $12$; $s^6{=}1$, $t^2{=}1{\neq}t$, $tst{=}s^{-1}$, $st{\neq}ts$ (diédrico).\\
\bottomrule
\end{tabular}
\caption{El Pilar~2 en \lean{NeoRiemannian.lean} / \lean{SixThirty.lean}, sin \lean{sorry} y con auditoría
de axiomas $[\mathrm{propext},\mathrm{Classical.choice},\mathrm{Quot.sound}]$. El dual de 6-30 es el único
punto donde el trabajo va más allá de la formalización; su alcance y la relación con el linaje de Fiore se
enuncian en la Sección~\ref{sec:novelty}.}
\label{tab:p2}
\end{table}
     % \S\ref{sec:pillar2}
% pillar3_scales_tuning_es.tex -- cuerpo del Pilar 3 (español), en la línea de casa.
% Sin preámbulo: usa las macros del maestro. La honestidad vive en la Sección \ref{sec:novelty}.
% Término: "máxima uniformidad" = maximal evenness (glosado una vez).

\section{Pilar 3: escalas y afinación}\label{sec:pillar3}

\question{P\textsubscript{0}. ¿Por qué es especial la escala diatónica, y es eso un solo hecho o varios que
casualmente coinciden?}

Cuatro condiciones de regularidad sobre una escala ---estar generada por un único intervalo, ser bien
formada, tener la propiedad de Myhill, ser de máxima uniformidad--- seleccionan la diatónica, y la última
reconecta con la autocorrelación del Pilar~1. Este pilar verifica por máquina cada una, nombra la instancia
allí donde sólo se demuestra una instancia, y termina en los retículos de afinación donde el hilo de la
dualidad del Pilar~2 regresa.

\subsection{Escalas generadas y bien formadas}\label{sec:p3-genwf}

Apilar $k$ quintas justas desde una altura de partida y reducir módulo $12$ (\lean{stack}) genera la
diatónica en $k=7$ (\lean{diatonic\_generated\_by\_fifth}) y la pentatónica en $k=5$
(\lean{pentatonic\_generated\_by\_fifth}). Una escala es \emph{bien formada} cuando tiene exactamente dos
tamaños de paso (\lean{IsWellFormed}). La diatónica lo es, con pasos $\{1,2\}$ semitonos
(\lean{diatonic\_isWellFormed}, \lean{diatonic\_stepSizes}), y también la pentatónica, con pasos $\{2,3\}$
(\lean{pentatonic\_isWellFormed}); apilar demasiado lo rompe --- seis quintas dan tres tamaños de paso
(\lean{sixFifths\_not\_wellFormed}). Que una escala generada tenga a lo sumo \emph{tres} tamaños de paso es
el teorema de los tres huecos (Steinhaus), que verificamos en $\Ztwelve$ (\lean{threeGap\_chromatic};
Sós~\cite{Steinhaus}). La teoría de buena formación de Carey--Clampitt~\cite{CareyClampitt} es el marco.

\subsection{La propiedad de Myhill}\label{sec:p3-myhill}

\question{P\textsubscript{1}. La buena formación cuenta tamaños de paso. La propiedad de Myhill cuenta algo
más fino ---y separa la diatónica de su prima uniforme.}

Una escala tiene la propiedad de Myhill cuando todo intervalo genérico no nulo (segunda, tercera, \dots)
aparece en exactamente dos tamaños específicos (\lean{HasMyhill}). La diatónica la tiene
(\lean{diatonic\_hasMyhill}); la escala de tonos enteros \emph{no} (\lean{wholeTone\_not\_hasMyhill}), pues
su único paso de tono colapsa la variedad. Por Carey--Clampitt, Myhill equivale a la buena formación no
degenerada; formalizamos la instancia diatónica y el contraejemplo de tonos enteros, y señalamos la
equivalencia general (que necesita la maquinaria de palabras de Christoffel/tres huecos, escasa en Mathlib)
como aún no formalizada.

\subsection{Máxima uniformidad}\label{sec:p3-me}

\question{P\textsubscript{2}. ``Repartir $k$ notas lo más uniformemente posible'': ¿qué energía lo hace
preciso, y singulariza de verdad la diatónica?}

La máxima uniformidad (\emph{maximal evenness}, Clough--Douthett~\cite{CloughDouthett}) es la más profunda
de las cuatro condiciones, y la formalizamos en tres capas, siendo explícitos sobre qué energía hace el
trabajo\,\audio{modes_messiaen.wav}. \emph{(i) El predicado.} Una escala es de máxima uniformidad cuando
cada intervalo genérico toma a lo sumo dos tamaños específicos y son enteros consecutivos
(\lean{IsMaxEven}); la diatónica y la escala de tonos enteros lo cumplen (\lean{diatonic\_isMaxEven},
\lean{wholeTone\_isMaxEven}), de modo que la máxima uniformidad contiene estrictamente a Myhill (la escala
de tonos enteros es de máxima uniformidad pero, por la~\S\ref{sec:p3-myhill}, no es de Myhill). \emph{(ii)
La energía de huecos-de-paso.} Para un potencial estrictamente convexo $V$, la energía
$\sum_i V(\text{hueco}_i)$ sobre los huecos de paso circulares se minimiza exactamente en las
configuraciones \emph{balanceadas}, $\max-\min\le1$ (\lean{isMinimizer\_iff\_balanced}), vía un lema de
intercambio por convexidad estricta y un descenso bien fundado. Esto caracteriza el multiconjunto de pasos
balanceado ---condición \emph{necesaria} para la máxima uniformidad, pero no suficiente, y lo decimos en la
Sección~\ref{sec:novelty}. \emph{(iii) La energía de todos-los-pares.} La energía que sí selecciona la
diatónica de forma única es el hamiltoniano de todos-los-pares de Douthett--Krantz
$E(V,A)=\sum_{a<b}V(\mathrm{cdist}(a,b))$ sobre todas las distancias por pares. Para todo $V$ convexo
decreciente la diatónica lo minimiza, y para todo $V$ estrictamente tal es el minimizador \emph{único}
salvo transposición e inversión (\lean{diatonic\_ground\_state}, \lean{diatonic\_unique}) ---un solo
enunciado que vale a la vez para todos los potenciales admisibles, la $V$-independencia. El motor es una
identidad de suma por partes doble (\lean{abel\_bridge}): dos $k$-subconjuntos tienen iguales totales de
vector de intervalos $\binom{k}{2}$, así que su diferencia suma cero y el salto de energía colapsa a una
suma de sumas doblemente acumuladas contra la convexidad de $V$. El mismo teorema vale para las demás
cardinalidades de máxima uniformidad en $\Ztwelve$ ---pentatónica ($k=5$), tonos enteros ($k=6$),
octatónica ($k=8$), que se unen a la diatónica ($k=7$)--- (\lean{pentatonic\_unique},
\lean{wholetone\_unique}, \lean{octatonic\_unique}), vía un núcleo independiente de la cardinalidad.
(Equivalentemente, la diatónica maximiza de forma única $|\Ahat(5)|$, Quinn~\cite{Quinn} ---el hilo de la
autocorrelación del Pilar~1, que regresa.)

\subsection{Teoría de temperamentos regulares: vals y retículos de comas}\label{sec:p3-rtt}

\question{P\textsubscript{3}. Doce quintas casi cierran la octava; el hueco es la coma pitagórica. ¿Cuál es
el álgebra de ese ``casi''?}

Modela una altura por sus exponentes en los generadores primos: un val es un mapa $\mathbb{Z}$-lineal que
manda los generadores a cuentas de paso (\lean{v12}, \lean{v5}, \lean{v7}). Doce quintas menos siete
octavas es la coma pitagórica $\mathrm{pc}$ (\lean{closure\_defect\_eq\_pc}); el temperamento igual de doce
tonos es exactamente el val que la anula, y la coma genera todo el núcleo de rango $1$,
$\ker v_{12}=\mathbb{Z}\cdot\mathrm{pc}$ (\lean{ker\_v12\_eq\_span\_pc}), siendo el val de $12$ el único que
la mata (\lean{only\_twelve\_tempers\_pc}). En el límite de $5$ entra la coma sintónica: el temperamento
igual de $12$ la tempera (\lean{v12\_5\_tempers\_syntonic}), cuatro quintas igualan una tercera mayor salvo
esa coma (\lean{meantone\_defect\_eq\_syntonic}) ---así que el temperamento igual de $12$ es mesotónico---
y el núcleo de comas es de rango $2$, generado por las comas sintónica y pitagórica
(\lean{ker\_v12\_5\_eq\_span}). Reducir la quinta módulo $12$ cierra el círculo de quintas sobre las doce
clases de altura (\lean{circle\_of\_fifths\_complete}, Figura~\ref{fig:p3-fifths}), y cuatro quintas se
reducen a la tercera mayor (\lean{four\_fifths\_eq\_major\_third}). El emparejamiento val--coma ---un
retículo generado y su anulador--- es el eco en $\mathbb{Z}$-módulos de la dualidad de centralizador del
Pilar~2, la rima de la~\S\ref{sec:p2-thread}.

\begin{figure}[ht]
\centering
\begin{tikzpicture}[scale=1.2,every node/.style={font=\small}]
  \draw[gray!55] (0,0) circle (1);
  \foreach \i in {0,...,11}{
    \pgfmathtruncatemacro{\pc}{mod(7*\i,12)}
    \coordinate (n\i) at ({90-30*\i}:1);
    \filldraw[ctA] (n\i) circle (0.035);
    \node at ({90-30*\i}:1.18) {\pc};}
  \foreach \i in {0,...,11}{
    \draw[ctA,->,>=Stealth,thin] ({90-30*\i}:0.92) arc ({90-30*\i}:{90-30*\i-26}:0.92);}
\end{tikzpicture}
\caption{El círculo de quintas: avanzar por la quinta reducida ($7$ semitonos) visita las doce clases de
altura (\lean{circle\_of\_fifths\_complete}); las etiquetas de los nodos son las clases de altura
alcanzadas en orden de quintas. El cierre es el núcleo de comas de rango $1$ de
\lean{ker\_v12\_eq\_span\_pc}.}
\label{fig:p3-fifths}
\end{figure}

\begin{table}[ht]
\centering\scriptsize\setlength{\tabcolsep}{4pt}
\rowcolors{2}{ctA!7}{white}
\begin{tabular}{@{}p{0.40\textwidth} p{0.52\textwidth}@{}}
\toprule
\rowcolor{ctA!22}
\textbf{Teorema} & \textbf{Enunciado}\\
\midrule
\lean{diatonic\_generated\_by\_fifth}, \lean{pentatonic\_generated\_by\_fifth} & diatónica / pentatónica generadas por quintas apiladas.\\
\lean{diatonic\_isWellFormed}, \lean{sixFifths\_not\_wellFormed} & bien formada $=$ dos tamaños de paso; el límite.\\
\lean{threeGap\_chromatic} & teorema de los tres huecos / Steinhaus en $\Ztwelve$.\\
\lean{diatonic\_hasMyhill}, \lean{wholeTone\_not\_hasMyhill} & propiedad de Myhill: diatónica sí, tonos enteros no.\\
\lean{diatonic\_isMaxEven}, \lean{wholeTone\_isMaxEven} & máxima uniformidad (predicado); máx.\ unif.\ $\supsetneq$ Myhill.\\
\lean{isMinimizer\_iff\_balanced} & minimizador de energía de huecos $\iff$ multiconjunto balanceado (necesario, no suficiente).\\
\lean{diatonic\_ground\_state}, \lean{diatonic\_unique} & energía de todos-los-pares: diatónica el minimizador único salvo $T/I$, $\forall$ $V$ admisible.\\
\lean{pentatonic\_unique}, \lean{wholetone\_unique}, \lean{octatonic\_unique} & lo mismo, en $k\in\{5,6,8\}$, vía \lean{abel\_bridge}.\\
\lean{closure\_defect\_eq\_pc}, \lean{ker\_v12\_eq\_span\_pc} & coma pitagórica $=$ defecto de cierre; núcleo de rango $1$.\\
\lean{meantone\_defect\_eq\_syntonic}, \lean{ker\_v12\_5\_eq\_span} & coma sintónica / mesotónico; núcleo de rango $2$ en el límite de $5$.\\
\lean{circle\_of\_fifths\_complete}, \lean{four\_fifths\_eq\_major\_third} & doce quintas generan $\Ztwelve$; $4$ quintas $=$ tercera mayor.\\
\bottomrule
\end{tabular}
\caption{El Pilar~3 en \lean{DiatonicScale.lean}, \lean{MaximalEvenness.lean},
\lean{AllPairsEvenness.lean}, \lean{Temperament.lean}, \lean{CircleOfFifths.lean}, sin \lean{sorry};
auditoría de axiomas $[\mathrm{propext},\mathrm{Classical.choice},\mathrm{Quot.sound}]$, salvo los censos
finitos de todos-los-pares, que cargan además $\mathrm{Lean.ofReduceBool}$ (\lean{native\_decide}),
señalado en la Sección~\ref{sec:novelty}.}
\label{tab:p3}
\end{table}
 % \S\ref{sec:pillar3}

\section{Novedad, alcance, y lo que este trabajo no afirma}\label{sec:novelty}

\question{¿Qué hay aquí de matemática nueva, qué es primera formalización de matemática conocida, y qué es
mera organización? Lo respondemos en un solo sitio, sin concesiones, para que el encuadre no se desdibuje.}

\paragraph{La contribución, dicha sin rodeos.} Este es un trabajo de formalización. Con las dos
excepciones de abajo, cada resultado es una demostración verificada por máquina de un teorema
publicado, citado a su fuente en su pilar y recogido en la Tabla~\ref{tab:ledger}. La biblioteca está sin
\lean{sorry}; la auditoría de axiomas es $[\mathrm{propext},\,\mathrm{Classical.choice},\,
\mathrm{Quot.sound}]$ en todo, salvo los censos finitos de máxima uniformidad, que cargan además
$\mathrm{Lean.ofReduceBool}$ de \lean{native\_decide} ---señalado aquí como en el cuerpo, con el motor
analítico que hay detrás (una identidad de suma por partes doble) mantenido axioma-limpio.

\paragraph{Lo genuinamente nuestro (1): la auto-dualidad de 6-30.} El único punto donde creemos ir más
allá de la formalización es la auto-dualidad de orden $12$ centrada en el tritono de la clase de Petrushka
$6$-$30=\{0,1,3,6,7,9\}$ (\S\ref{sec:p2-sixthirty}). No es inversionalmente simétrica pero la fija el
tritono $T_6$; su órbita $T/I$ colapsa por tanto a $12$, el tritono cae en el núcleo de la acción, y el
centralizador de la acción reducida es diédrico de orden $12$. No conocemos una caracterización explícita
previa de este dual, ni un dual contextual de un acorde simétrico verificado por máquina. Lo enunciamos
\emph{con mesura}, como en la Nota~\#1~\cite{Note1}: es \emph{complementario} al linaje de Fiore, no algo que se les escapara.
Fiore--Satyendra~\cite{FS} y Fiore--Noll~\cite{FN11,FN25} esquivan cada uno el régimen de órbita reducida
con centro en el tritono; \cite{FN11} hasta tabula la misma clase de conjuntos como un punto y estudia su
recubrimiento de orden $2$, sin formar nunca el dual de orden $12$ del hexacordo-como-punto.
Berry--Fiore~\cite{BF} demuestran la misma dualidad de centralizador para ciclos hexatónicos ---otro
objeto (estabilizador aumentado, dual $S_3$ de orden $6$). El único riesgo de colisión que queda, los
grupos conmutantes generalizados de Peck~\cite{Peck}, lo dejamos matizado, no descartado; los grupos
conmutantes octatónicos publicados en este linaje son de orden $8$, no $12$. Lo que afirmamos es la
caracterización explícita y su dual verificado por máquina, no la prioridad sobre la teoría general.

\paragraph{Lo genuinamente nuestro (2): la taxonomía de fase unificada.} Cada ingrediente de la taxonomía
de fase es clásico (Patterson, Lewin, Babbitt, Amiot). Lo que creemos inédito \emph{como un solo
enunciado} es la unificación: clavar cada invariante a la porción exacta de $\Ahat$ que retiene y
certificar las dos fronteras de la escalera ---el colapso ciego a la magnitud (homometría) y la cúspide del
invariante completo--- una junto a la otra, con una sola auditoría de axiomas (\S\ref{sec:p1-capstone}).
Esto es organización, no matemática nueva; no reclamamos prioridad sobre ninguno de los resultados de base.

\paragraph{Un punto de encuadre menor.} En la teoría de escalas dejamos explícito que la caracterización
por energía convexa de la máxima uniformidad se parte en una forma de huecos-de-paso degenerada y una forma
de todos-los-pares que sí selecciona la diatónica (\S\ref{sec:p3-me}); es una aclaración de qué objeto
hace el trabajo, no un teorema.

\paragraph{Trabajo previo en asistentes de demostración.} La única formalización musical en Lean que
conocemos (Prismriver~\cite{Prismriver}) aborda la acción $T/I$; las tres capas de aquí ---los invariantes
de Fourier/vector de intervalos, la dualidad neo-riemanniana, y la teoría de escalas y afinación--- son,
que sepamos, primeras formalizaciones. Un lema de teoría de grupos reutilizable, el caso regular del
teorema del centralizador de Wielandt (\lean{centralizing\_fixedPoint\_eq\_one}), no estaba en Mathlib y se
aporta aquí.

\begin{table}[H]
\centering\footnotesize\setlength{\tabcolsep}{5pt}
\rowcolors{2}{ctA!7}{white}
\begin{tabular}{@{}p{0.37\textwidth} p{0.31\textwidth} p{0.17\textwidth}@{}}
\toprule
\rowcolor{ctA!22}
\textbf{Resultado} & \textbf{Fuente clásica} & \textbf{Estatus}\\
\midrule
\multicolumn{3}{@{}l}{\emph{Pilar 1 --- invariantes de Fourier}}\\
Vector de intervalos $=\Fi|\Ahat|^2$; notas comunes & Amiot~\cite{Amiot}; Lewin~\cite{Lewin},Rahn~\cite{Rahn} & 1.\textsuperscript{a} formaliz.\\
Teorema del hexacordo de Babbitt & Babbitt; Lewin~\cite{Lewin1960} & 1.\textsuperscript{a} formaliz.\\
Escalas profundas; rareza rítmica & Gamer/Browne~\cite{Browne}; Chemillier--Truchet~\cite{ChemillierTruchet} & 1.\textsuperscript{a} formaliz.\\
Homometría $=$ problema de la fase & Patterson~\cite{Patterson}; Forte~\cite{Forte} & 1.\textsuperscript{a} formaliz.\\
Taxonomía de fase unificada (un enunciado) & --- (este trabajo) & \textbf{nuestro: organización}\\
\multicolumn{3}{@{}l}{\emph{Pilar 2 --- dualidad transformacional}}\\
Rep.\ regular de $T/I$; voice-leading de $PLR$ & Lewin~\cite{Lewin}; Cohn~\cite{Cohn} & 1.\textsuperscript{a} formaliz.\\
Lema del caso regular de Wielandt & Dixon--Mortimer~\cite{DM} & 1.\textsuperscript{a} formaliz.$^\dagger$\\
Dualidad CFS: $PLR=$ centralizador$(T/I)$ & Crans--Fiore--Satyendra~\cite{CFS} & 1.\textsuperscript{a} formaliz.\\
\textbf{Auto-dualidad de 6-30 con centro en el tritono} & --- (este trabajo) & \textbf{nuestro: resultado nuevo}\\
\multicolumn{3}{@{}l}{\emph{Pilar 3 --- escalas y afinación}}\\
Myhill; bien formada; tres huecos & Carey--Clampitt~\cite{CareyClampitt}; tres huecos & 1.\textsuperscript{a} formaliz.\\
Evenness máxima (predicado; energía) & Clough--Douthett~\cite{CloughDouthett}; Douthett--Krantz~\cite{DouthettKrantz} & 1.\textsuperscript{a} formaliz.\\
\quad distinción huecos-de-paso vs todos-los-pares & --- (este trabajo) & nuestro: organización\\
Retículos de comas (pitagórica, sintónica) & teoría de temperamentos regulares & 1.\textsuperscript{a} formaliz.\\
\bottomrule
\end{tabular}
\caption{El registro de honestidad. Cada fila es una primera formalización de matemática clásica, citada a
su fuente, salvo las dos filas \textbf{nuestro} (un resultado nuevo, una de organización) y el punto de
encuadre menor. $^\dagger$ El lema del caso regular de Wielandt no estaba en Mathlib; aportarlo es una
pequeña contribución de formalización reutilizable más allá del contenido musical.}
\label{tab:ledger}
\end{table}

\section{La frontera abierta}\label{sec:frontier}

\question{¿Qué podemos enunciar pero aún no certificar en el núcleo, y hacia dónde apunta cada hilo?}

En el espíritu del programa más amplio, tratamos la frontera de lo formalizado como un resultado, no como
una carencia que ocultar. Quedan abiertas cuatro direcciones.

\begin{itemize}\itemsep2pt
  \item \textbf{La familia auto-dual de $n$ general.} La caracterización de dualidad reducida tras $6$-$30$
        es independiente del universo; $6$-$30$ es la instancia no trivial mínima de una familia con dual
        $\cong\Dgrp{n/2}$ en el universo $n$ (el conteo es OEIS A032239). Un enunciado uniforme en Lean, y
        el isomorfismo con nombre $\Dgrp{12}/Z\cong\Dgrp{6}$ (hoy una laguna de Mathlib, atestiguada en
        GAP), siguen abiertos.
  \item \textbf{Evenness máxima en $(k,n)$ general.} La unicidad $V$-independiente de todos-los-pares está
        certificada sobre $\Ztwelve$ para $k\in\{5,6,7,8\}$ por censo finito; el motor analítico (la
        identidad de suma por partes doble) es general, pero el teorema en $(k,n)$ general necesita el
        argumento abstracto de mayorización, y sustituir los censos por una prueba de núcleo retiraría la
        única dependencia de \lean{Lean.ofReduceBool}.
  \item \textbf{Fase y auto-dualidad.} La simetría $T_6$ es exactamente soporte en frecuencias pares de
        $\Ahat$, un hecho de magnitud; pero la auto-dualidad necesita información de fase que $|\Ahat|$ no
        puede ver ---la misma ceguera que la relación $Z$. Quedan abiertos un relato de la auto-dualidad en
        el espacio de fase, y un teorema en Lean para la completitud de clase de conjuntos (hoy
        computacional).
  \item \textbf{Del análisis a la composición.} La biblioteca es un motor de análisis; el generador
        maximalmente uniforme, el diagrama de fases de las escalas bien formadas, y la navegación por el grupo
        dual de un acorde simétrico son constructivos y apuntan a un programa de síntesis aguas abajo.
\end{itemize}

\section*{Disponibilidad de datos y código}
Las fuentes de Lean, los guiones testigo de Sage/GAP y las figuras acompañan a este trabajo; cada teorema
de Lean es comprobable con \lean{\#print axioms} y cada guion es re-ejecutable. Las tres notas que lo
componen están archivadas en Zenodo (\cite{Note1,Note2,Note3}).

\begin{thebibliography}{99}
% --- el hilo DFT / Fourier ---
\bibitem{Loeffler} D. Loeffler, Mathlib \lean{ZMod.dft} (Apache-2.0) --- the discrete Fourier
  transform on \lean{ZMod N} as a \lean{LinearEquiv}.
\bibitem{Amiot} E. Amiot, \emph{Music Through Fourier Space: Discrete Fourier Transform in Music
  Theory}, Computational Music Science, Springer, 2016.
\bibitem{Patterson} A. L. Patterson, ``A Fourier Series Method for the Determination of the
  Components of Interatomic Distances in Crystals,'' \emph{Phys. Rev.} \textbf{46} (1934),
  372--376.
\bibitem{MAAA} J. Mandereau, D. Ghisi, E. Amiot, M. Andreatta, C. Agon, ``Z-relation and
  homometry in musical distributions,'' \emph{J. Math. \& Music} \textbf{5} (2011), no.~2,
  83--98.
\bibitem{Quinn} I. Quinn, ``General equal-tempered harmony,'' \emph{Perspectives of New Music}
  \textbf{44}--\textbf{45} (2006--2007).
% --- teoría de conjuntos / invariantes ---
\bibitem{Forte} A. Forte, \emph{The Structure of Atonal Music}, Yale Univ. Press, 1973.
\bibitem{Lewin} D. Lewin, \emph{Generalized Musical Intervals and Transformations}, Yale Univ.
  Press, 1987 (reimpr.\ Oxford Univ. Press, 2007); GIS $=$ acción simplemente transitiva, p.~157.
\bibitem{Lewin1959} D. Lewin, ``Intervallic relations between two collections of notes,''
  \emph{J. Music Theory} \textbf{3} (1959), no.~2, 298--301.
\bibitem{Lewin1960} D. Lewin, ``Re: The intervallic content of a collection of notes\dots,''
  \emph{J. Music Theory} \textbf{4} (1960), no.~1, 98--101.
\bibitem{Rahn} J. Rahn, \emph{Basic Atonal Theory}, Longman, 1980.
\bibitem{Browne} R. Browne, ``Tonal implications of the diatonic set,'' \emph{In Theory Only}
  \textbf{5} (1981), no.~6--7, 3--21; C. Gamer, ``Some combinational resources of equal-tempered
  systems,'' \emph{J. Music Theory} \textbf{11} (1967), 32--59.
\bibitem{ChemillierTruchet} M. Chemillier, C. Truchet, ``Computing the rhythmic oddity property,''
  \emph{Inform. Process. Lett.} \textbf{86} (2003), no.~5, 255--261; G. T. Toussaint, \emph{The
  Geometry of Musical Rhythm}, CRC Press, 2013.
% --- transformacional / dualidad ---
\bibitem{Cohn} R. Cohn, ``Neo-Riemannian operations, parsimonious trichords, and their Tonnetz
  representations,'' \emph{J. Music Theory} \textbf{41} (1997), no.~1, 1--66.
\bibitem{CFS} A. S. Crans, T. M. Fiore, R. Satyendra, ``Musical actions of dihedral groups,''
  \emph{Amer. Math. Monthly} \textbf{116} (2009), no.~6, 479--495 (arXiv:0711.1873).
\bibitem{DM} J. D. Dixon, B. Mortimer, \emph{Permutation Groups}, Springer GTM 163, 1996
  (\S4.2, Theorem 4.2A).
\bibitem{Cam} P. J. Cameron, \emph{Permutation Groups}, LMS Student Texts 45, CUP, 1999.
\bibitem{BF} C. Berry, T. M. Fiore, ``Hexatonic systems and dual groups in mathematical music
  theory,'' \emph{Involve} \textbf{11} (2018), no.~2, 253--270 (arXiv:1602.02577).
\bibitem{Peck} R. W. Peck, ``Generalized commuting groups,'' \emph{J. Music Theory} \textbf{54}
  (2010), no.~2, 143--177.
\bibitem{FS} T. M. Fiore, R. Satyendra, ``Generalized contextual groups,'' \emph{Music Theory
  Online} \textbf{11}.3 (2005).
\bibitem{FN11} T. M. Fiore, T. Noll, ``Commuting groups and the topos of triads,'' in
  \emph{Mathematics and Computation in Music} (MCM 2011), LNCS \textbf{6726}, Springer, 2011,
  pp.~69--83 (arXiv:1102.1496).
\bibitem{FN25} T. M. Fiore, T. Noll, \emph{et al.}, ``Transformations of triads and seventh
  chords: group extensions and duality,'' arXiv:2507.20811 (2025).
% --- escalas / afinación ---
\bibitem{CareyClampitt} N. Carey, D. Clampitt, ``Aspects of well-formed scales,'' \emph{Music
  Theory Spectrum} \textbf{11} (1989), no.~2, 187--206.
\bibitem{CloughDouthett} J. Clough, J. Douthett, ``Maximally even sets,'' \emph{J. Music Theory}
  \textbf{35} (1991), no.~1/2, 93--173.
\bibitem{DouthettKrantz} J. Douthett, R. Krantz, ``Maximally even sets and configurations: common
  threads in mathematics, physics, and music,'' \emph{J. Comb. Optim.} \textbf{14} (2007), no.~4,
  385--410.
\bibitem{BakBruinsma} P. Bak, R. Bruinsma, ``One-dimensional Ising model and the complete devil's
  staircase,'' \emph{Phys. Rev. Lett.} \textbf{49} (1982), 249--251.
\bibitem{Steinhaus} V. T. S\'os, ``On the distribution mod $1$ of the sequence $n\alpha$,''
  \emph{Ann. Univ. Sci. Budapest. E\"otv\"os} \textbf{1} (1958), 127--134; S. \'Swierczkowski,
  ``On successive settings of an arc on the circumference of a circle,'' \emph{Fund. Math.}
  \textbf{46} (1959), 187--189 (teorema de los tres huecos / Steinhaus).
\bibitem{Erlich} P. Erlich, ``A middle path between just intonation and the equal temperaments,''
  \emph{Xenharmonik\^on} \textbf{18} (2006); teoría de temperamentos regulares (vals y comas).
% --- asistentes de demostración / trabajo previo ---
\bibitem{Prismriver} L. Aniva, C. Wang, \emph{Prismriver: formalization of music theory and
  algorithmic composition in Lean~4}, arXiv:2606.19936 (2026).
% --- nuestras notas ---
\bibitem{Note1} C. Mar\'in, ``A tritone-centered self-duality: the set class 6-30 and its
  order-12 neo-Riemannian dual,'' Nota~\#1 (serie music-math), 2026
  (DOI:10.5281/zenodo.20820962).
\bibitem{Note2} C. Mar\'in, ``The phase taxonomy of pitch-class set invariants,'' Nota~\#2
  (serie music-math), 2026 (DOI:10.5281/zenodo.20826774).
\bibitem{Note3} C. Mar\'in, ``Fourier spectra and pitch symmetry'' (el espectro del Tonnetz),
  Nota~\#3 (serie music-math), 2026.
\end{thebibliography}

\end{document}
